\subsection{Построение геометрии}

Для построения фигуры используется правильный тетраэдр с центром в начале координат. Его четыре вершины выбраны симметрично, поэтому расстояния между любыми двумя исходными вершинами одинаковы. На каждом ребре отмечаются две точки: на расстоянии одной трети длины ребра от каждого из его концов.

В программе исходные вершины имеют координаты
\(A_0=(1,1,1)\), \(A_1=(1,-1,-1)\),
\(A_2=(-1,1,-1)\) и \(A_3=(-1,-1,1)\).
Новая вершина на направленном ребре от \(A_i\) к \(A_j\) вычисляется как
\begin{equation}
  P_{ij}=\frac{2A_i+A_j}{3}, \qquad i\ne j.
  \label{eq:truncation}
\end{equation}
Здесь \(i,j\in\{0,1,2,3\}\). Для каждого ребра используются оба направления,
поэтому получаются 12 различных точек. Длина исходного ребра равна
\(2\sqrt{2}\), а длина любого ребра усечённой фигуры — \(2\sqrt{2}/3\).

Вершины исходного тетраэдра отсекаются плоскостями, проходящими через отмеченные точки. Каждая из четырёх отсечённых вершин образует новую треугольную грань. Каждая исходная треугольная грань после удаления трёх углов превращается в шестиугольник. Усечение именно на одной трети ребра обеспечивает одинаковую длину сторон треугольников и шестиугольников.

Полученная модель содержит 12 вершин, 18 рёбер и 8 граней: четыре правильных треугольника и четыре правильных шестиугольника. Фигура является выпуклой. Каждое ребро принадлежит двум граням, а в каждой вершине сходятся три ребра.

Геометрия хранится отдельно от интерфейса в файлах \path{tetrahedron.hpp} и \path{tetrahedron.cpp}. Модель представлена массивом пространственных вершин, списком граней и списком уникальных рёбер. Для каждой грани задаётся нормаль, направленная наружу. Вершины грани упорядочиваются по кругу относительно её центра, что обеспечивает последовательный обход контура и корректную заливку.
